% ==============================================================================
% Fichier : 04_referentiels.tex
% Description : Normes et Référentiels Usuels en Audit des SI
% ==============================================================================

\chapter{Normes et Référentiels Usuels en Audit des Systèmes d'Information}\label{chap:04_referentiels}

L'auditeur a besoin de \og{}briques\fg{} techniques et organisationnelles pour évaluer la maturité
d'un SI. Les référentiels fournissent les bonnes pratiques reconnues internationalement et
constituent le vocabulaire commun entre l'auditeur, la direction et les équipes IT.
Ce chapitre présente d'abord les trois piliers fondateurs, COSO, COBIT et la série ISO~27000,
puis neuf référentiels complémentaires qui enrichissent la palette de l'auditeur selon les
domaines couverts~: valeur des investissements, risques, maturité des processus, gestion des
services et méthodologie d'audit.

% ==============================================================================
\section{COSO~: Le Modèle de Contrôle Interne}
\label{sec:coso}
% ==============================================================================

\subsection{Présentation}

\begin{boxdef}
\textbf{COSO (Committee of Sponsoring Organizations)~:} Cadre de référence conceptuel du
contrôle interne, publié en~1992 (révisé en~2013). C'est le modèle standard pour la gouvernance
d'entreprise, souvent exigé par les lois sur la gouvernance (ex~: Sarbanes-Oxley aux USA).
\end{boxdef}

Bien que généraliste, COSO est le socle sur lequel s'appuient les autres référentiels
informatiques.

\subsection{Les Trois Objectifs}

Le contrôle interne selon COSO vise à donner une assurance raisonnable sur la réalisation
de trois objectifs~:
\begin{enumerate}
    \item \textbf{Réalisation et Optimisation des opérations} (Performance).
    \item \textbf{Fiabilité des informations financières} (Comptabilité).
    \item \textbf{Conformité aux lois et règlements} (Juridique).
\end{enumerate}

\subsection{Les Cinq Composantes (Le Cube COSO)}

Pour atteindre ces objectifs, l'organisation doit mettre en place 5~composantes liées~:

\begin{table}[h!]
\centering
\small
\begin{tabular}{p{4cm}p{8cm}}
    \toprule
    \textbf{Composante} & \textbf{Description} \\
    \midrule
    \textbf{Environnement de Contrôle} & La culture, l'éthique, la compétence du personnel.
        (Le \og{}ton donné par le sommet\fg{}). \\
    \textbf{Évaluation des Risques} & Identification et analyse des risques qui pourraient
        empêcher d'atteindre les objectifs. \\
    \textbf{Activités de Contrôle} & Les politiques et procédures concrètes (Séparation des
        tâches, autorisations). \\
    \textbf{Information et Communication} & La circulation de l'information nécessaire au
        pilotage (Qualité des données). \\
    \textbf{Pilotage (Monitoring)} & Surveillance continue du système de contrôle interne
        (Audits, Supervision). \\
    \bottomrule
\end{tabular}
\caption{Les 5~composantes du référentiel COSO}
\end{table}

\begin{boxlizbiz}
\textbf{Application Lizbiz's~:}
L'audit du système de gestion des commandes chez Lizbiz's vérifiera si les
5~composantes COSO sont présentes~: existe-t-il une procédure (Activité de contrôle)~?
Les risques de fraude sont-ils évalués (Évaluation des risques)~? Les managers supervisent-ils
les anomalies (Pilotage)~?
\end{boxlizbiz}

% ==============================================================================
\section{COBIT~: La Gouvernance des SI}
\label{sec:cobit}
% ==============================================================================

\subsection{Présentation}

\begin{boxdef}
\textbf{COBIT (Control Objectives for Information and Related Technologies)~:} Référentiel
de gouvernance et de management des SI développé par l'ISACA. Il fait le pont entre les
métiers et l'IT.
\end{boxdef}

C'est le référentiel de choix pour l'auditeur SI car il propose une liste d'objectifs de
contrôle précis pour chaque processus informatique.

\subsection{Objectifs et Apports}

COBIT répond à la question~: \og{}Comment aligner l'IT sur la stratégie de l'entreprise~?\fg{}
Il structure le SI en domaines et processus, offrant à l'auditeur une grille de lecture
universelle.

\subsection{Structure de COBIT (Version 5 / 2019)}

COBIT divise la gouvernance et le management en 5~domaines principaux~:

\begin{table}[h!]
\centering
\small
\begin{tabular}{p{2.8cm}p{10.2cm}}
    \toprule
    \textbf{Domaine} & \textbf{Contenu (Exemples de processus)} \\
    \midrule
    \textbf{Evaluer, Diriger et Surveiller (EDM)} & Gouvernance stratégique
        (EDM01~: Gouvernance du cadre I\&T). \\
    \midrule
    \textbf{Planifier et Organiser (PO)} & Stratégie, Architecture, Budget
        (PO01~: Gérer le cadre de management IT). \\
    \midrule
    \textbf{Acquérir et Implémenter (AI)} & Projets, Développement, Intégration
        (AI06~: Gérer les changements). \\
    \midrule
    \textbf{Délivrer et Servir (DS)} & Exploitation, Support, Sécurité
        (DS01~: Gérer les opérations). \\
    \midrule
    \textbf{Surveiller, Évaluer et Analyser (MEA)} & Performance, Conformité
        (MEA02~: Surveiller et évaluer le contrôle interne). \\
    \bottomrule
\end{tabular}
\caption{Les domaines de COBIT}
\end{table}
\subsection{Utilisation par l'Auditeur}

Pour chaque processus (ex~: DS5 \og{}Gérer la sécurité\fg{}), COBIT fournit~:
\begin{itemize}
    \item Un objectif de contrôle détaillé.
    \item Des indicateurs de performance (KPI).
    \item Un modèle de maturité (de~0~: Inexistant à~5~: Optimisé).
\end{itemize}

\subsection{La Cascade d'Objectifs COBIT}

L'un des points forts de COBIT est sa capacité à lier la stratégie métier aux objectifs
techniques~:
\begin{enumerate}
    \item \textbf{Besoins des parties prenantes} (Bénéfices, Risques, Ressources).
    \item \textbf{Objectifs d'entreprise} (ex~: Portefeuille de produits et services structurés).
    \item \textbf{Objectifs d'alignement} (ex~: Qualité des informations financières liées à l'IT).
\end{enumerate}

% ==============================================================================
\section{La Série ISO~27000~: La Sécurité Normative}
\label{sec:iso27000}
% ==============================================================================

Alors que COBIT gère l'IT, la série ISO~27000 se concentre spécifiquement sur la
\textbf{Sécurité de l'Information}. C'est une famille de normes internationales dont
les membres sont indissociables~: ISO~27001 pose les exigences, ISO~27002 fournit
les mesures pour les satisfaire, ISO~27005 fournit la méthode pour en justifier
la sélection par le risque, et ISO~27006 encadre les organismes qui délivrent
la certification.

% ------------------------------------------------------------------------------
\subsection{Vue d'Ensemble}
\label{subsec:iso27000-vue}
% ------------------------------------------------------------------------------

\begin{table}[h!]
\centering
\small
\renewcommand{\arraystretch}{1.4}
\begin{tabular}{|p{2.4cm}|p{3.2cm}|p{6.4cm}|}
    \hline
    \rowcolor{blue!10}
    \textbf{Norme} & \textbf{Rôle} & \textbf{Contenu principal} \\ \hline
    \textbf{ISO~27001} & Exigences (certifiable) &
        Spécifie les exigences du SMSI~: politique, gestion des risques, objectifs,
        mesures (Annexe~A), amélioration continue. \\ \hline
    \textbf{ISO~27002} & Code de bonnes pratiques &
        93~mesures de sécurité détaillées avec lignes directrices d'implémentation
        (compagnon opérationnel d'ISO~27001). \\ \hline
    \textbf{ISO~27005} & Gestion des risques &
        Cadre méthodologique pour identifier, analyser, évaluer et traiter les risques
        de sécurité de l'information. \\ \hline
    \textbf{ISO~27006} & Certification &
        Exigences pour les organismes qui auditent et certifient les SMSI. \\ \hline
\end{tabular}
\caption{Les normes clés de la famille ISO~27000}
\end{table}

% ------------------------------------------------------------------------------
\subsection{ISO~27001. Exigences du Système de Management de la Sécurité de l'Information}
\label{subsec:iso27001}
% ------------------------------------------------------------------------------

\begin{boxdef}
\textbf{ISO~27001~:} Norme certifiable qui définit les \emph{exigences} d'un Système
de Management de la Sécurité de l'Information (SMSI). Elle impose à l'organisation de
mettre en place un cycle structuré~: définir le contexte, évaluer les risques, sélectionner
des mesures de traitement, les implémenter, les surveiller et les améliorer en continu
(roue de Deming PDCA).
\end{boxdef}

L'Annexe~A d'ISO~27001 liste les 93~mesures de contrôle issues d'ISO~27002 parmi
lesquelles l'organisation sélectionne celles applicables à son contexte, en justifiant
les exclusions dans la \textbf{Déclaration d'Applicabilité (SoA)}.

\paragraph{Utilisation par l'auditeur.}
ISO~27001 fournit les \emph{critères d'audit}~: l'auditeur vérifie que chaque exigence
de la norme est satisfaite. La SoA est le document pivot~: elle indique quelles mesures
sont retenues, pourquoi, et quel est leur état d'implémentation. C'est sur elle que
repose le \emph{gap analysis} de sécurité.

\begin{boxlizbiz}
\textbf{Projet Lizbiz's~:}
Suite à l'audit, la direction de Lizbiz's souhaite rassurer ses clients et décide
de se faire certifier \textbf{ISO~27001}. L'auditeur utilisera l'Annexe~A comme
grille d'évaluation pour le futur audit de certification, en construisant la SoA
avec les équipes sécurité.
\end{boxlizbiz}

% ------------------------------------------------------------------------------
\subsection{ISO~27002. Code de Bonnes Pratiques pour la Sécurité de l'Information}
\label{subsec:iso27002}
% ------------------------------------------------------------------------------

\begin{boxdef}
\textbf{ISO~27002:2022~:} Compagnon opérationnel d'ISO~27001. Là où ISO~27001 pose
les \emph{exigences}, ISO~27002 fournit les \emph{lignes directrices} concrètes pour
implémenter chacune des 93~mesures référencées en Annexe~A. Ce n'est pas une norme
certifiable~: c'est un code de bonnes pratiques.
\end{boxdef}

La version~2022 restructure profondément le référentiel (de 114~mesures en 14~domaines
vers 93~mesures en 4~thèmes) et introduit 11~nouvelles mesures~: surveillance des
activités, gestion des configurations, sécurité du cloud, préparation aux incidents.

\paragraph{Structure.} Les 93~mesures sont réparties en quatre thèmes~:
\begin{itemize}
    \item \textbf{Organisationnel} (37~mesures)~: politiques, rôles, gestion des risques,
        intelligence des menaces~;
    \item \textbf{Humain} (8~mesures)~: sélection, sensibilisation, clauses contractuelles~;
    \item \textbf{Physique} (14~mesures)~: périmètre, équipements, supports~;
    \item \textbf{Technologique} (34~mesures)~: authentification, chiffrement, journalisation,
        filtrage web, gestion des vulnérabilités.
\end{itemize}

\paragraph{Utilisation par l'auditeur.} ISO~27002 est la référence de premier niveau pour
tout audit de sécurité~: elle fournit la liste des mesures applicables une fois la SoA
établie, justifie chaque objectif de contrôle par une recommandation normative
internationale (renforçant l'opposabilité du rapport), et sert de base au
\emph{gap analysis}~: écart entre les mesures recommandées et celles effectivement
implémentées. Elle est structurellement compatible avec COBIT, permettant une
gouvernance IT unifiée.

% ------------------------------------------------------------------------------
\subsection{ISO~27005. Gestion des Risques Liés à la Sécurité de l'Information}
\label{subsec:iso27005}
% ------------------------------------------------------------------------------

\begin{boxdef}
\textbf{ISO~27005:2022~:} Norme méthodologique pour la gestion des risques de sécurité
de l'information, alignée sur ISO~31000 (management du risque en général). Elle boucle
le triptyque ISO~27001~/~ISO~27002~/~ISO~27005~: les risques identifiés selon ISO~27005
justifient le choix des mesures ISO~27002 retenues dans la SoA d'ISO~27001.
\end{boxdef}

ISO~27005 ne prescrit pas de méthode unique~: elle est compatible avec des approches
quantitatives (EBIOS~RM, MEHARI, OCTAVE) et qualitatives.

\paragraph{Structure.} Le processus ISO~27005 suit cinq étapes cycliques~:
\begin{enumerate}
    \item \textbf{Établissement du contexte~:} périmètre, critères d'évaluation,
        appétit au risque~;
    \item \textbf{Identification des risques~:} actifs informationnels, menaces,
        vulnérabilités, impacts potentiels~;
    \item \textbf{Analyse des risques~:} estimation de la vraisemblance et de l'impact,
        calcul du niveau de risque~;
    \item \textbf{Évaluation des risques~:} comparaison au seuil d'acceptabilité,
        priorisation~;
    \item \textbf{Traitement des risques~:} sélection des mesures (réduire, accepter,
        transférer, éviter), plan de traitement.
\end{enumerate}

\paragraph{Utilisation par l'auditeur.} L'auditeur mobilise ISO~27005 pour évaluer
la qualité du processus de gestion des risques de l'entité~: la méthodologie est-elle
formalisée et révisée~? La cartographie est-elle exhaustive~? Les risques résiduels
ont-ils été acceptés formellement par le niveau décisionnel approprié~? ISO~27005
est complémentaire de Risk~IT~: Risk~IT couvre les risques IT au sens large
(gouvernance, valeur, processus)~; ISO~27005 se concentre sur les seuls risques
de sécurité de l'information.

% ------------------------------------------------------------------------------
\subsection{Importance pour l'Audit}
\label{subsec:iso27000-audit}
% ------------------------------------------------------------------------------

Les trois normes opérationnelles de la famille forment un triptyque indissociable
du point de vue de l'auditeur~:

\begin{table}[h!]
\centering
\small
\renewcommand{\arraystretch}{1.4}
\begin{tabular}{|p{2.4cm}|p{4.0cm}|p{5.6cm}|}
    \hline
    \rowcolor{blue!10}
    \textbf{Norme} & \textbf{Rôle en audit} & \textbf{Question clé} \\ \hline
    ISO~27001 & Critères d'audit (exigences) &
        Les exigences du SMSI sont-elles toutes satisfaites~? \\ \hline
    ISO~27002 & Référentiel de mesures (recommandations) &
        Les mesures de sécurité retenues sont-elles implémentées et efficaces~? \\ \hline
    ISO~27005 & Fondement de la sélection des mesures &
        Le processus de gestion des risques est-il rigoureux et formalisé~? \\ \hline
\end{tabular}
\caption{Rôles respectifs des normes ISO~27000 en audit}
\end{table}
% ==============================================================================
\section{Référentiels Complémentaires}
\label{sec:referentiels-complementaires}
% ==============================================================================

Les référentiels présentés dans les sections précédentes forment un socle. Ils sont complétés par
un écosystème normatif spécialisé~: chacun des référentiels suivants aborde le SI sous un angle
précis, et leur valeur augmente lorsqu'ils sont utilisés de façon articulée plutôt qu'isolée.

% ------------------------------------------------------------------------------
\subsection{Val~IT. Gouvernance de la Valeur des Investissements IT}
\label{subsec:val-it}
% ------------------------------------------------------------------------------

\begin{boxdef}
\textbf{Val~IT~:} Référentiel publié par l'ISACA en~2006 (révisé en~2008), conçu comme
complément de COBIT. Là où COBIT gouverne les \emph{processus} IT, Val~IT gouverne la
\emph{valeur} que ces processus créent pour l'organisation.
\end{boxdef}

Val~IT répond à la question fondamentale des directions générales~: \emph{l'investissement
consenti dans l'informatique produit-il le retour attendu~?}

\paragraph{Structure.} Val~IT s'articule en trois domaines couvrant 22~processus de
gestion~: la \emph{gouvernance de la valeur} (cadre de décision), la \emph{gestion du
portefeuille} (priorisation des projets) et la \emph{gestion des bénéfices} (réalisation
effective des gains après mise en service). Le \textbf{business case} en est la pièce
maîtresse~: aucun projet ne devrait être lancé sans une estimation documentée de sa valeur
nette.

\paragraph{Utilisation par l'auditeur.} L'auditeur mobilise Val~IT pour évaluer~:
\begin{itemize}
    \item l'existence et la qualité des business cases des projets IT en cours ou livrés~;
    \item la gouvernance du portefeuille IT~: les projets sont-ils priorisés selon des critères
        explicites alignés sur la stratégie~?
    \item le suivi de la réalisation des bénéfices~: une fois l'application livrée, vérifie-t-on
        que les gains promis ont bien été obtenus~?
\end{itemize}

Val~IT s'appuie sur COBIT pour l'évaluation des processus sous-jacents et est complémentaire de
COSO~: COSO évalue le contrôle interne~; Val~IT évalue si les investissements qui l'alimentent
sont justifiés.

% ------------------------------------------------------------------------------
\subsection{Risk~IT. Gouvernance des Risques Liés aux Technologies de l'Information}
\label{subsec:risk-it}
% ------------------------------------------------------------------------------

\begin{boxdef}
\textbf{Risk~IT~:} Publié par l'ISACA en~2009, troisième pilier du triptyque
COBIT~/~Val~IT~/~Risk~IT. Il couvre l'ensemble des risques IT susceptibles d'affecter les
objectifs de l'organisation, qu'ils soient techniques, organisationnels ou stratégiques.
\end{boxdef}

Risk~IT distingue explicitement les \emph{risques IT propres} (panne, vulnérabilité,
obsolescence) des \emph{risques métier liés à l'IT} (dysfonctionnements de processus métier
causés par une défaillance IT).

\paragraph{Structure.} Risk~IT est organisé en trois domaines~:
\begin{itemize}
    \item \textbf{Gouvernance des risques IT (RG)~:} cadre de gestion des risques, appétit
        au risque, culture du risque~;
    \item \textbf{Évaluation des risques IT (RE)~:} identification, analyse et cartographie~;
    \item \textbf{Réponse aux risques IT (RR)~:} mise en œuvre des réponses (acceptation,
        réduction, transfert, évitement).
\end{itemize}

\paragraph{Utilisation par l'auditeur.} Risk~IT permet d'évaluer la maturité du processus de
gestion des risques de l'entité~: l'organisation dispose-t-elle d'une cartographie formalisée et
mise à jour~? L'appétit au risque est-il défini~? Les risques résiduels sont-ils acceptés
formellement par le niveau décisionnel approprié~?

\begin{boxlizbiz}
\textbf{Projet Lizbiz's~:}
La DSI de Lizbiz's ne dispose d'aucune cartographie des risques IT. L'auditeur s'appuiera sur le
domaine~RE de Risk~IT pour structurer la cartographie préliminaire de la mission, puis sur le
domaine~RG pour recommander l'instauration d'un appétit au risque formalisé.
\end{boxlizbiz}

% ------------------------------------------------------------------------------
\subsection{NIST SP~800-53. Catalogue de Contrôles de Sécurité}
\label{subsec:nist-800-53}
% ------------------------------------------------------------------------------

\begin{boxdef}
\textbf{NIST SP~800-53~:} Publication spéciale du National Institute of Standards and
Technology (États-Unis), révision~5 (2020). Catalogue exhaustif de contrôles de sécurité et
de protection de la vie privée, à portée initialement fédérale américaine mais devenu une
référence internationale.
\end{boxdef}

À la différence d'ISO~27001 qui définit des exigences organisationnelles, SP~800-53 spécifie des
contrôles \emph{techniques et procéduraux granulaires} portant directement sur le système.

\paragraph{Structure.} La révision~5 comprend 20~familles de contrôles~; parmi les plus
pertinentes pour l'audit~:
\begin{itemize}
    \item \textbf{AC} (\emph{Access Control})~: 25~contrôles sur la gestion des accès,
        des privilèges et des sessions~;
    \item \textbf{AU} (\emph{Audit and Accountability})~: 16~contrôles sur la journalisation
        et la traçabilité~;
    \item \textbf{CM} (\emph{Configuration Management})~: inventaires, lignes de base de
        configuration, gestion des changements~;
    \item \textbf{IR} (\emph{Incident Response})~: détection, réponse et récupération~;
    \item \textbf{SC} (\emph{System and Communications Protection})~: chiffrement, séparation
        des réseaux.
\end{itemize}

\paragraph{Utilisation par l'auditeur.} SP~800-53 sert de check-list lors des audits techniques~:
vérification de la configuration des systèmes, construction de la matrice de couverture des
contrôles, évaluation des accès dans des environnements complexes ou multi-couches. Il est souvent
utilisé en complément d'ISO~27001~: ISO~27001 fournit le cadre de management~;
SP~800-53 fournit le détail technique des contrôles.

% ------------------------------------------------------------------------------
\subsection{CMMi. Modèle d'Évaluation de la Maturité des Processus IT}
\label{subsec:cmmi}
% ------------------------------------------------------------------------------

\begin{boxdef}
\textbf{CMMI (\emph{Capability Maturity Model Integration})~:} Développé par le Software
Engineering Institute (SEI) de l'Université Carnegie~Mellon. La version actuelle,
CMMI~v2.0 (2018), mesure le \emph{degré de maturité} des processus d'une organisation sur
une échelle en cinq niveaux~: il ne prescrit pas de solutions techniques mais garantit que les
résultats sont prévisibles et améliorables.
\end{boxdef}

\paragraph{L'échelle de maturité.}
\begin{enumerate}
    \item \textbf{Initial~:} processus imprévisibles, réactifs, dépendants des individus.
    \item \textbf{Géré~:} processus gérés au niveau des projets (planification, surveillance).
    \item \textbf{Défini~:} processus documentés, standardisés, proactifs à l'échelle de
        l'organisation.
    \item \textbf{Géré quantitativement~:} processus mesurés à l'aide d'indicateurs
        quantitatifs.
    \item \textbf{En optimisation~:} amélioration continue basée sur les données~; innovation
        intégrée au processus.
\end{enumerate}

\paragraph{Utilisation par l'auditeur.} CMMi est l'outil de référence pour évaluer la maturité
de la fonction Études (qualité des processus de développement, tests, documentation) et pour
positionner la DSI sur l'échelle 1--5, donnant une vue synthétique communicable à la direction.
L'indice de maturité $M = \tau_c \times 5$ utilisé dans la feuille de synthèse de l'outil d'audit
reflète directement l'esprit de cette échelle. CMMi est complémentaire de COBIT (gouvernance des
processus) et de Val~IT (valeur des processus).

% ------------------------------------------------------------------------------
\subsection{ITIL. Bibliothèque des Bonnes Pratiques des Services IT}
\label{subsec:itil}
% ------------------------------------------------------------------------------

\begin{boxdef}
\textbf{ITIL (\emph{Information Technology Infrastructure Library})~:} Recueil de bonnes
pratiques en gestion des services IT (\emph{ITSM}), maintenu par AXELOS. La version
actuelle, \textbf{ITIL~4} (2019), abandonne la structure en cinq livres au profit d'un
\emph{Système de Valeur des Services} (SVS) articulé autour de 34~pratiques de gestion.
\end{boxdef}

ITIL structure la gestion des services IT du point de vue de leur \emph{valeur pour
l'utilisateur}. Il est la référence mondiale des centres de services, des équipes d'exploitation
et des helpdesks. La distinction avec COBIT est nette~: COBIT se demande \emph{ce que} l'IT
doit accomplir~; ITIL se demande \emph{comment} le faire dans une logique de service.

\paragraph{Pratiques clés pertinentes pour l'audit.}
\begin{itemize}
    \item \textbf{Gestion des incidents~:} processus de restauration rapide du service~; référence
        pour l'audit des procédures d'incident~;
    \item \textbf{Gestion des problèmes~:} identification et résolution des causes profondes des
        incidents récurrents~;
    \item \textbf{Gestion des changements~:} contrôle de toute modification apportée au SI~;
    \item \textbf{Gestion des niveaux de service~:} définition, suivi et révision des SLA~;
    \item \textbf{Gestion des actifs IT~:} inventaire et suivi du cycle de vie~;
    \item \textbf{Gestion des connaissances~:} documentation et capitalisation des savoirs IT.
\end{itemize}

\paragraph{Utilisation par l'auditeur.} L'auditeur utilise ITIL comme grille de référence pour
évaluer la qualité des processus d'exploitation en régime normal, vérifier que les SLA sont
définis et respectés, et auditer la gestion des changements~: toute modification de production
devrait suivre un processus de demande, d'évaluation, d'autorisation, de test et de traçabilité.

% ------------------------------------------------------------------------------
\subsection{ISO~20000. Système de Management des Services IT}
\label{subsec:iso20000}
% ------------------------------------------------------------------------------

\begin{boxdef}
\textbf{ISO/IEC~20000-1:2018~:} Norme internationale de gestion des services IT, directement
dérivée de BS~15000. ISO~20000 est à ITIL ce qu'ISO~27001 est à ISO~27002~: elle transforme les
bonnes pratiques en exigences \emph{certifiables}.
\end{boxdef}

ISO~20000 couvre la planification, la conception, la transition, la livraison et l'amélioration
des services IT en imposant une approche systémique et documentée.

\paragraph{Utilisation par l'auditeur.}
\begin{itemize}
    \item \emph{Audit de certification~:} vérification de la conformité du SMS aux exigences de
        la norme en vue d'une certification tierce partie~;
    \item \emph{Audit de conformité interne~:} évaluation de l'adéquation des processus de
        gestion de services (helpdesk, incidents, changements, niveaux de service) aux exigences
        normatives~;
    \item \emph{Benchmark de maturité~:} comparaison du niveau de maturité du SMS avec les
        exigences normatives, sans ambition de certification.
\end{itemize}

Une organisation certifiée ISO~20000 offre à l'auditeur une présomption de maturité des
processus de gestion de services~: les certifications ISO~20000 et les rapports SOC~1
constituent des éléments probants importants lors de l'audit de fournisseurs de services IT
externalisés.

% ------------------------------------------------------------------------------
\subsection{ITAF. Cadre de Référence pour l'Audit des Systèmes d'Information}
\label{subsec:itaf}
% ------------------------------------------------------------------------------

\begin{boxdef}
\textbf{ITAF (\emph{IT Assurance Framework})~:} Publié par l'ISACA, cadre méthodologique
spécifique à la \emph{pratique de l'audit des SI}. Là où COBIT, ITIL ou ISO~27001 gouvernent
l'IT, l'ITAF gouverne la \emph{façon dont l'auditeur} conduit ses missions. Il complète les
normes générales d'audit (IIA, INTOSAI) par des directives propres aux systèmes d'information.
\end{boxdef}

\paragraph{Structure.} L'ITAF est organisé en trois niveaux~:
\begin{itemize}
    \item \textbf{Standards généraux~:} principes d'indépendance, compétences requises, éthique
        professionnelle~;
    \item \textbf{Standards de performance~:} planification, supervision, documentation,
        collecte de preuves~;
    \item \textbf{Standards de reporting~:} structure du rapport, conclusions,
        recommandations.
\end{itemize}

\paragraph{Utilisation par l'auditeur.} L'ITAF est la référence professionnelle directe de
l'auditeur certifié CISA. Il définit le niveau de diligence attendu pour chaque phase de la
mission, précise les types de preuves acceptables (observation, inspection de documents,
ré-exécution de contrôles, interrogation) et encadre la communication des conclusions (clarté,
objectivité, complétude, opportunité). Il est également la référence pour l'évaluation des
\emph{contrôles compensatoires}~: lorsqu'un contrôle attendu est absent, l'auditeur recherche
si un contrôle compensatoire produit une assurance équivalente. L'ITAF est méthodologiquement
neutre~: il est compatible avec tous les référentiels techniques (COBIT, ISO, NIST).

% ==============================================================================
\section{Synthèse~: Comment Choisir~?}
\label{sec:synthese-referentiels}
% ==============================================================================

L'auditeur combine ces référentiels selon le périmètre et les objectifs de la mission. Les trois
piliers fondateurs couvrent les angles stratégiques essentiels~:

\begin{table}[h!]
\centering
\begin{tabular}{p{3cm}p{4cm}p{5cm}}
    \toprule
    \textbf{Référentiel} & \textbf{Angle d'attaque} & \textbf{Question clé pour l'auditeur} \\
    \midrule
    \textbf{COSO} & Gouvernance Globale & Les contrôles internes sont-ils bien conçus~? \\
    \textbf{COBIT} & Processus IT & L'informatique est-elle bien gérée~? \\
    \textbf{ISO~27001} & Sécurité Info & Le système est-il sécurisé et certifiable~? \\
    \bottomrule
\end{tabular}
\caption{Comparatif des référentiels fondateurs}
\end{table}

Les référentiels complémentaires affinent cette lecture selon le domaine investigué~:

\begin{table}[H]
\centering
\small
\renewcommand{\arraystretch}{1.4}
\caption{Synthèse des référentiels complémentaires~: portée, usage en audit et domaine associé}
\label{tab:synthese-referentiels}
\begin{tabular}{|p{2.4cm}|p{2.8cm}|p{4.2cm}|p{2.8cm}|}
    \hline
    \rowcolor{gray!20}
    \textbf{Référentiel} & \textbf{Portée principale} &
    \textbf{Usage en audit} & \textbf{Domaine associé} \\ \hline
    Val~IT & Valeur des investissements IT &
    Business cases, portefeuille IT, bénéfices réalisés & Gouvernance DSI \\ \hline
    Risk~IT & Risques IT (gouvernance) &
    Cartographie des risques IT, appétit au risque & Sécurité / Risques \\ \hline
    NIST SP~800-53 & Contrôles de sécurité (technique) &
    Check-list contrôles techniques, matrice de couverture & Sécurité \\ \hline
    ISO~27002 & Bonnes pratiques sécurité &
    Gap analysis mesures de sécurité, critères d'audit & Sécurité \\ \hline
    ISO~27005 & Risques sécurité de l'information &
    Qualité du processus de gestion des risques sécurité & Sécurité \\ \hline
    CMMi & Maturité des processus IT &
    Maturité DSI/Étude, indice de maturité 0--5 & Études / Synthèse \\ \hline
    ITIL & Gestion des services IT &
    Exploitation, incidents, changements, SLA & Support / Continuité \\ \hline
    ISO~20000 & SMS certifiable (ITSM) &
    Audit de certification/conformité ITSM, SOC prestataires & Support / Cloud \\ \hline
    ITAF & Méthodologie d'audit des SI &
    Diligences, preuves, reporting, contrôles compensatoires & Toutes missions \\ \hline
\end{tabular}
\end{table}

Le tableau synoptique ci-dessous présente le comparatif complet des référentiels fondateurs en
fonction de leur cible, de leur focus auditeur et du livrable type attendu~:

\begin{table}[h!]
\centering
\small
\renewcommand{\arraystretch}{1.5}
\begin{tabular}{|p{2.5cm}|p{2.5cm}|p{4.5cm}|p{3.5cm}|}
    \hline
    \rowcolor{blue!10} \textbf{Référentiel} & \textbf{Cible} & \textbf{Focus Auditeur} &
    \textbf{Livrable type} \\ \hline
    \textbf{COSO} & DG / Audit Interne & Contrôle Interne Global &
    Rapport de conformité Loi Sarbanes-Oxley \\ \hline
    \textbf{COBIT} & DSI / Métiers & Gouvernance et Alignement &
    Matrice de maturité des processus SI \\ \hline
    \textbf{ISO~27001} & RSSI / Sécurité & Management de la Sécurité &
    Certification SMSI (Audit de conformité) \\ \hline
\end{tabular}
\caption{Tableau comparatif des référentiels d'audit fondateurs}
\end{table}

\begin{boxlizbiz}
\textbf{Synthèse pour Lizbiz's~:}
Un audit complet du SI de Lizbiz's mobilisera~: COSO pour la revue du contrôle interne global,
COBIT pour la gouvernance IT et la maturité des processus, ISO~27001/27002 pour la sécurité,
Val~IT pour évaluer la pertinence des investissements informatiques récents, Risk~IT et
ISO~27005 pour la cartographie des risques, ITIL pour l'audit de l'exploitation et des niveaux
de service, CMMi pour positionner la maturité de la DSI, et ITAF comme boussole méthodologique
tout au long de la mission.
\end{boxlizbiz}
